
\paragraph{n-n}\label{n-n}

\[
\begin{aligned}
\bm{A}_{\Psi} &= \int \bm{\Psi}^{\mathrm{t}}\bm{\Psi}\,\mathrm{d}A %\\
&= \int \begin{bmatrix}
 1 & 0 & 0 \\
 0 & \psi_{yy}^2 + \psi_{zy}^2 & \psi_{yy}\psi_{zy} + \psi_{zz}\psi_{yz} \\
 0 & \psi_{yy}\psi_{zy} + \psi_{zz}\psi_{yz}  & \psi_{zz}^2 + \psi_{yz}^2
 \end{bmatrix} \, \mathrm{d} A %\\
&= \begin{bmatrix}
A & 0 & 0 \\
0 & A_y & A_{yz} \\
0 & A_{zy} & A_z
\end{bmatrix}
\end{aligned}
\] The \textbf{three} independent constants \(A, A_y\) and \(A_z\) are
of interest. If \(\bm{\Psi} = \bm{1}\) then all constants reduce to the
cross sectional area \(A\).

\paragraph{m-m}\label{m-m}

\[
\begin{aligned}
\bm{I} &= \SectionIntegral{
     [\bm{r}^\times][\bm{r}^\times]^{\mathrm{t}} 
}
= \SectionIntegral{ \left[\begin{matrix}
y^{2} + z^{2} & 0 & 0\\
0 & z^{2} & - z y \\
0 & - z y & \,\,y^{2}
\end{matrix}\right]
}
% &= \int (\bm{r}\cdot\bm{r})\bm{1} - \bm{r}\otimes \bm{r}\, \mathrm{d}A \\
% &=\left[\delta_{\alpha \beta} \mathbf{1}-\FrameBasis_\alpha \otimes \FrameBasis_\beta\right]\int_{\Omega} \xi_\alpha \xi_\beta \, \mathrm{d}A  \\
\\
&\triangleq I_0 \FrameBasis \otimes \FrameBasis + I_{\alpha \beta} \FrameBasis_\alpha \otimes \FrameBasis_\beta  \\
\end{aligned}
\]
\textbf{three} independent constants:
\[
\begin{aligned}
I_0    &\triangleq  ...\\
I_y    &\triangleq  I_{22} = \int z^2\, \mathrm{d} A \\
I_z    &\triangleq  I_{33} = \int y^2\, \mathrm{d} A \\
I_{zy} &\triangleq -I_{23} = \int yz \, \mathrm{d} A
\end{aligned}
\]
Note
\[
\bm{I} = \SectionIntegral{\bm{r}\cdot\bm{r}\bm{1} - \bm{r}\otimes\bm{r}}
\]

\paragraph{w-w}\label{w-w}

\[
\bm{C} = \int \bm{\Phi}^{\mathrm{t}} \bm{\Phi} \, \mathrm{d} A
= \int \begin{bmatrix}
\varphi^2 & 0 \\
0 & \varphi_{, y}^2 + \varphi_{, z}^2
\end{bmatrix} \, \mathrm{d} A
= \begin{bmatrix}
C_{\varphi} & 0 \\
0 & C_{\alpha}
\end{bmatrix}
\] where \textbf{two} independent constants \(C_{\varphi},C_{\alpha}\)
are defined.

\begin{itemize}
\tightlist
\item
  If \(\bm{\varphi} = \tilde{\varphi} \,\FrameBasis\) then \(C_{\varphi}\) is the warping
  constant of \cite[p39]{vlasov1961thinwalled} 
  \[
  C_{\varphi} = \SectionIntegral{ \tilde{\varphi}^2 }
  \] 
\item \textcite{gruttmann2000theory} refer to \(C_{\alpha}\) with the notation ``\(I_0 + I_T\)''
  \[
  C_{\alpha}  = \SectionIntegral{ \tilde{\varphi}_{, y}^2 + \tilde{\varphi}_{, z}^2} = J_S - J
  \]
\end{itemize}

\paragraph{n-m}\label{n-m}

\[
\bm{Q}_{\Psi}
= \SectionIntegral{ - \bm{\Psi}^{\mathrm{t}} [\bm{r}^\times] }
= \SectionIntegral{ \begin{bmatrix}
0 & z & -y \\
- z\psi_{yy} + y \psi_{zy} & 0 & 0 \\
-z \psi_{yz} + y \psi_{zz} & 0 & 0
\end{bmatrix} \, \mathrm{d} A
 = \begin{bmatrix}
0 & Q_{xy} & Q_{xz} \\
Q_{yx} & 0 & 0 \\
Q_{zx} & 0 & 0
\end{bmatrix}
}
\] 
where \textbf{four} independent constants
\(Q_{xy}, Q_{xz}, Q_{yx}, Q_{zx}\) have been defined.

\begin{itemize}
\item
  If \(\bm{\Psi} = \bm{1}\) then \(\bm{Q} = A\hat{\bm{r}}^{\times}\) \[
  Q_{yx} = \int -z\,\mathrm{d} A
  \]
\item
  If, additionally, the reference curve passes through the centroid as
  in \cite{simo1991geometricallyexact} then \(\bm{Q} = 0\)
\end{itemize}

\paragraph{n-w}\label{n-w}

\[
\bm{R} = \SectionIntegral{ \bm{\Psi}^\mathrm{t}\bm{\Phi} 
}
= \SectionIntegral{\begin{bmatrix}
\tilde{\varphi} & 0 \\
0 & \psi_{yy} \varphi_{, y} + \psi_{zy} \varphi_{, z} \\
0 & \psi_{yz} \varphi_{, y} + \psi_{zz} \varphi_{, z}
\end{bmatrix}
}
= \begin{bmatrix}
R_{\varphi} & 0 \\
0 & R_{y} \\
0 & R_{z}
\end{bmatrix}
\] 
where \textbf{three} independent constants have been defined.

\begin{itemize}
\item
  If \(\bm{\Psi} = \bm{1}\) and \(\varphi\) is referred to the \emph{shear
  center} (e.g., \cite{gruttmann1998geometrical}) 
  \begin{equation*}
  \begin{aligned}
  R_{\varphi} &= 0 \\
  R_y &= \int \tilde{\varphi}_{,y} = -(\tilde{z} - \hat{z})A \\
  R_z &= \int \tilde{\varphi}_{,z} =  (\tilde{y} - \hat{y})A
  \end{aligned}
  \end{equation*}
\item
  If \(\bm{\Psi} = \bm{1}\) and \(\varphi\) is referred to the \emph{centroid}
  (i.e.~\(\varphi = \hat{\varphi}\) like \cite{gruttmann2000theory}) then:
  \[
  R_y = \int \hat{\varphi}_{,y} =    \hat{z}A \\
  R_z = \int \hat{\varphi}_{,z} =  - \hat{y}A
  \] where \(\hat{z}\) and \(\hat{y}\) are the coordinates of the
  centroid.
\end{itemize}

\paragraph{m-w}\label{m-w}

\[
\bm{S} 
= \SectionIntegral{ \bm{r}^{\times} \bm{\Phi} } 
= \SectionIntegral{ \begin{bmatrix}
0 & -z \varphi_{, y} + y \varphi_{, z} \\
 z \varphi & 0 \\
-y \varphi & 0
\end{bmatrix}
}
= \begin{bmatrix}
0 & S_{\alpha} \\
S_{y} & 0 \\
S_{z} & 0
\end{bmatrix}
\] 
introducing \textbf{three} constants \(S_{\alpha}, S_y\) and \(S_z\).

\begin{itemize}
\tightlist
\item
  If \(\varphi\) is referred to the centroid 
  \[
  \begin{aligned}
  S_y &= \int z \hat{\varphi} \, \mathrm{d} A =   \hat{I}_{z} \hat{z} - \hat{I}_{yz} \hat{y} \\
  S_z &= \int y \hat{\varphi} \, \mathrm{d} A = - \hat{I}_{y} \hat{y} + \hat{I}_{yz} \hat{z}
  \end{aligned}
  \]
\end{itemize}
